% GENERATED FILE. Edit canonical lessons, then run npm run generate:books.
% canonical-source-hash: fnv1a64:d5c7a1790f463d60
% canonical-lessons: JA-W13-tsu, JA-R13-tsu-contrast, JA-C13-chichi, JA-C13-haha, JA-C13-ani, JA-C13-ane, JA-C13-otouto, JA-C13-imouto, JA-C13-otto, JA-C13-tsuma, JA-C13-kodomo, JA-R13-family-a, JA-R13-family-b, JA-R13-family-c, JA-R13-family-d, JA-C13-family-reception, JA-C13-family-check

\chapter{Nine Family Words Through One New Kana}
\label{ch:ja-family-nine}
\hlchaptermodality{13}

\begin{chapteropening}
\textbf{By the end of this chapter:} \emph{I can hear, say, read, and write nine immediate-family words, with listening, speaking, reading, and writing scored independently.}

Nine official Irodori Starter family words built from earned hiragana plus one separately traced, copied, and recalled full-size \ja{つ}.
\end{chapteropening}

\section[tsu]{\ja{つ} — one full-size curve}

\label{lesson:JA-W13-tsu}



\begin{quote}
Write \textbf{\ja{ありがとう}}, name its older \emph{arigatashi} family, and point to the dakuten. Then write small \textbf{\ja{っ}} and retrieve one word from the body map.
\end{quote}

\subsection*{Script — notice and trace}
\begin{quote}
\textbf{\ja{つ}} — \emph{tsu}
\end{quote}

This is the full-size sign. Trace its one stroke: begin at the upper left, sweep right, turn downward, and finish toward the lower left.

\subsection*{Writing — copy, then recall}
Copy \textbf{\ja{つ}} once. Hide it for ten seconds, write it once, then place \textbf{\ja{つ}} beside small \textbf{\ja{っ}} and check that the full-size sign is clearly larger.

\begin{tcolorbox}[breakable,colback=teal!4,colframe=teal!35!black,title={Before you move on}]
Read it once: \emph{tsu}.

Source: \href{https://www.unicode.org/charts/PDF/U3040.pdf}{Unicode Hiragana chart}; stroke order per \emph{Hitsujun Shido no Tebiki} (Japanese Ministry of Education, 1958).
\end{tcolorbox}
\section[Practice]{Full \ja{つ}, small \ja{っ}}

\label{lesson:JA-R13-tsu-contrast}



\begin{quote}
Write distant \textbf{\ja{さ}} once from memory.
\end{quote}

\subsection*{Your turn — one contrast}
Write distant \textbf{\ja{さ}}, then \textbf{\ja{つ・っ}}. Say \emph{tsu} for the full-size sign; hold one silent beat for the small sign.

\begin{tcolorbox}[breakable,colback=teal!4,colframe=teal!35!black,title={Before you move on}]
Stop after one accurate pair.
\end{tcolorbox}
\section[chichi]{\ja{ちち} — my father}

\label{lesson:JA-C13-chichi}



\begin{quote}
Write known \textbf{\ja{ち}} once and give it one even mora.
\end{quote}

\subsection*{You'll Want to Know — hear it before reading it}
Hear \emph{chi-chi}. Point to “my father,” then say the two even morae once.

\subsection*{Your turn — read, copy, recall}
Read \textbf{\ja{ち} | \ja{ち}}. Copy \textbf{\ja{ちち}}, hide it, and write it once from the meaning.

\begin{tcolorbox}[breakable,colback=teal!4,colframe=teal!35!black,title={Before you move on}]
Source: \href{https://www.irodori.jpf.go.jp/assets/data/wordlist\_X.pdf}{Japan Foundation Irodori Starter word list, Lesson 4}.
\end{tcolorbox}
\section[haha]{\ja{はは} — my mother}

\label{lesson:JA-C13-haha}



\begin{quote}
Say and write \textbf{\ja{ちち}} once before meeting the paired family word.
\end{quote}

\subsection*{You'll Want to Know — hear and say}
Write distant \textbf{\ja{す}}. Hear \emph{ha-ha}, say it once, then contrast \emph{chichi}.

\subsection*{Your turn — read, copy, recall}
Read \textbf{\ja{は} | \ja{は}}. Copy \textbf{\ja{はは}}, hide it, and write it once from the meaning.

\begin{tcolorbox}[breakable,colback=teal!4,colframe=teal!35!black,title={Before you move on}]
Source: \href{https://www.irodori.jpf.go.jp/assets/data/wordlist\_X.pdf}{Japan Foundation Irodori Starter word list, Lesson 4}.
\end{tcolorbox}
\section[ani]{\ja{あに} — my elder brother}

\label{lesson:JA-C13-ani}



\begin{quote}
Retrieve \textbf{\ja{はは}} once from its meaning.
\end{quote}

\subsection*{You'll Want to Know — hear and say}
Thank a stranger with \textbf{\ja{ありがとうございます}}, name the polite-register choice, and retrieve any two body words. Then hear \emph{a-ni} and say both morae once.

\subsection*{Your turn — read, copy, recall}
Read \textbf{\ja{あ} | \ja{に}}. Copy \textbf{\ja{あに}}, hide it, and write it from the meaning.

\begin{tcolorbox}[breakable,colback=teal!4,colframe=teal!35!black,title={Before you move on}]
Source: \href{https://www.irodori.jpf.go.jp/assets/data/wordlist\_X.pdf}{Japan Foundation Irodori Starter word list, Lesson 4}.
\end{tcolorbox}
\section[ane]{\ja{あね} — my elder sister}

\label{lesson:JA-C13-ane}



\begin{quote}
Retrieve \textbf{\ja{あに}} once from its meaning.
\end{quote}

\subsection*{You'll Want to Know — hear and say}
Read the distant block \textbf{\ja{ございます}} once. Hear \emph{a-ne}, say it, then contrast \emph{a-ni}, my elder brother.

\subsection*{Your turn — read, copy, recall}
Read \textbf{\ja{あ} | \ja{ね}}. Copy \textbf{\ja{あね}}, hide it, and write it from the meaning.

\begin{tcolorbox}[breakable,colback=teal!4,colframe=teal!35!black,title={Before you move on}]
Source: \href{https://www.irodori.jpf.go.jp/assets/data/wordlist\_X.pdf}{Japan Foundation Irodori Starter word list, Lesson 4}.
\end{tcolorbox}
\section[otooto]{\ja{おとうと} — my younger brother}

\label{lesson:JA-C13-otouto}



\begin{quote}
Retrieve \textbf{\ja{あね}} once from its meaning.
\end{quote}

\subsection*{You'll Want to Know — hear it first}
Write distant \textbf{\ja{日}} once. Hear \emph{o-to-o-to}: o | to | o | to, with the long o occupying two beats.

\subsection*{Your turn — read, copy, recall}
Read \textbf{\ja{お} | \ja{と} | \ja{う} | \ja{と}}. Copy \textbf{\ja{おとうと}}, hide it, and write it once.

\begin{tcolorbox}[breakable,colback=teal!4,colframe=teal!35!black,title={Before you move on}]
Source: \href{https://www.irodori.jpf.go.jp/assets/data/wordlist\_X.pdf}{Japan Foundation Irodori Starter word list, Lesson 4}.
\end{tcolorbox}
\section[imooto]{\ja{いもうと} — my younger sister}

\label{lesson:JA-C13-imouto}



\begin{quote}
Say \textbf{\ja{おとうと}} once and keep its four even morae.
\end{quote}

\subsection*{You'll Want to Know — hear and contrast}
Write distant \textbf{\ja{本}} once. Hear \emph{i-mo-o-to}, say it, then contrast \emph{o-to-o-to}.

\subsection*{Your turn — read, copy, recall}
Read \textbf{\ja{い} | \ja{も} | \ja{う} | \ja{と}}. Copy \textbf{\ja{いもうと}}, hide it, and write it once.

\begin{tcolorbox}[breakable,colback=teal!4,colframe=teal!35!black,title={Before you move on}]
Source: \href{https://www.irodori.jpf.go.jp/assets/data/wordlist\_X.pdf}{Japan Foundation Irodori Starter word list, Lesson 4}.
\end{tcolorbox}
\section[otto]{\ja{おっと} — my husband}

\label{lesson:JA-C13-otto}



\begin{quote}
Retrieve \textbf{\ja{いもうと}} once from its meaning.
\end{quote}

\subsection*{You'll Want to Know — hear the held beat}
Write distant \textbf{\ja{言}} once. Hear \emph{o | hold | to} and say it without inserting a vowel into the held beat.

\subsection*{Your turn — read, copy, recall}
Read \textbf{\ja{お} | \ja{っ} | \ja{と}}. Copy \textbf{\ja{おっと}}, hide it, and write it once.

\begin{tcolorbox}[breakable,colback=teal!4,colframe=teal!35!black,title={Before you move on}]
Source: \href{https://www.irodori.jpf.go.jp/assets/data/wordlist\_X.pdf}{Japan Foundation Irodori Starter word list, Lesson 4}.
\end{tcolorbox}
\section[tsuma]{\ja{つま} — my wife}

\label{lesson:JA-C13-tsuma}



\begin{quote}
Retrieve \textbf{\ja{おっと}} once and notice the small \textbf{\ja{っ}}.
\end{quote}

\subsection*{You'll Want to Know — hear and say}
Write distant \textbf{\ja{五}} once. Hear \emph{tsu-ma} and say two morae, with a full vowel after \emph{ts}.

\subsection*{Your turn — read, copy, recall}
Read \textbf{\ja{つ} | \ja{ま}}. Copy \textbf{\ja{つま}}, hide it, and write it once.

\begin{tcolorbox}[breakable,colback=teal!4,colframe=teal!35!black,title={Before you move on}]
Source: \href{https://www.irodori.jpf.go.jp/assets/data/wordlist\_X.pdf}{Japan Foundation Irodori Starter word list, Lesson 4}.
\end{tcolorbox}
\section[kodomo]{\ja{こども} — child}

\label{lesson:JA-C13-kodomo}



\begin{quote}
Retrieve \textbf{\ja{つま}} once from its meaning.
\end{quote}

\subsection*{You'll Want to Know — hear and say}
Write distant \textbf{\ja{口}} once. Hear \emph{ko-do-mo}, choose “child,” then say the three even morae.

\subsection*{Your turn — read, copy, recall}
Read \textbf{\ja{こ} | \ja{ど} | \ja{も}}. Copy \textbf{\ja{こども}}, hide it, and write it once.

\begin{tcolorbox}[breakable,colback=teal!4,colframe=teal!35!black,title={Before you move on}]
Source: \href{https://www.irodori.jpf.go.jp/assets/data/wordlist\_X.pdf}{Japan Foundation Irodori Starter word list, Lesson 4}.
\end{tcolorbox}
\section[Practice]{Family retrieval A}

\label{lesson:JA-R13-family-a}



\begin{quote}
Say \textbf{\ja{おとうと}} once without looking back.
\end{quote}

\subsection*{Your turn — two directions}
Build distant \textbf{\ja{語}} once. Hear \emph{otooto} and point to “younger brother.” Read \textbf{\ja{こども}}, then write both family words from memory.

\begin{tcolorbox}[breakable,colback=teal!4,colframe=teal!35!black,title={Before you move on}]
Stop after one accurate pair.
\end{tcolorbox}
\section[Practice]{Family retrieval B}

\label{lesson:JA-R13-family-b}



\begin{quote}
Say \textbf{\ja{いもうと}} once without looking back.
\end{quote}

\subsection*{Your turn — three cues}
Read distant \textbf{\ja{日本語}} as \emph{nihongo} and recall that borrowed kanji can carry word-specific readings. Then hear “younger sister,” say \textbf{\ja{いもうと}}, read \textbf{\ja{ちち・はは}}, and write all three.

\begin{tcolorbox}[breakable,colback=teal!4,colframe=teal!35!black,title={Before you move on}]
Stop after one accurate set.
\end{tcolorbox}
\section[Practice]{Family retrieval C}

\label{lesson:JA-R13-family-c}



\begin{quote}
Say \textbf{\ja{おっと}} once without looking back.
\end{quote}

\subsection*{Your turn — three cues}
Write distant katakana \textbf{\ja{コ}}. Hear \emph{otto} and mark its held beat. Read \textbf{\ja{あに・あね}}, say each meaning, then write all three.

\begin{tcolorbox}[breakable,colback=teal!4,colframe=teal!35!black,title={Before you move on}]
Stop after one accurate set.
\end{tcolorbox}
\section[Practice]{Family retrieval D}

\label{lesson:JA-R13-family-d}



\begin{quote}
Say \textbf{\ja{つま}} once without looking back.
\end{quote}

\subsection*{Your turn — size carries sound}
Add distant \textbf{\ja{ー}} after \textbf{\ja{コ}} and state that it holds the vowel for one more beat. Then hear, read, and write \textbf{\ja{つま・おっと}} with a clear size contrast.

\begin{tcolorbox}[breakable,colback=teal!4,colframe=teal!35!black,title={Before you move on}]
Stop after one accurate pair.
\end{tcolorbox}
\section[Practice]{Nine family words — reception check}

\label{lesson:JA-C13-family-reception}



\begin{quote}
Retrieve \textbf{\ja{こども}} once from the English cue “child.”
\end{quote}

\subsection*{Your turn — listening}
Use sound only. Score listening before revealing any written word.

\subsection*{Your turn — reading}
Read without romanization. A failed section is repeated on its own.

\begin{tcolorbox}[breakable,colback=teal!4,colframe=teal!35!black,title={Before you move on}]
Keep listening and reading scores separate; repair only the weaker channel.
\end{tcolorbox}
\section[Practice]{Nine family words — production check}

\label{lesson:JA-C13-family-check}



\begin{quote}
Write distant katakana \textbf{\ja{ヒ}}, then read \textbf{\ja{コーヒー}} and recall why a loanword uses katakana. Hear, say, read, and write the seven later body words: \textbf{\ja{あたま・かみ・は・かた・おなか・せなか・こし}}.
\end{quote}

\subsection*{Your turn — speaking}
Use no written answer while speaking. Check task completion and mora timing.

\subsection*{Writing}
No section compensates for another. Repeat only the section below its own pass mark.

\begin{tcolorbox}[breakable,colback=teal!4,colframe=teal!35!black,title={Before you move on}]
Keep speaking and writing scores separate; repair only the weaker channel.

Source: \href{https://www.irodori.jpf.go.jp/assets/data/wordlist\_X.pdf}{Japan Foundation Irodori Starter word list, Lesson 4}.
\end{tcolorbox}
